%----------------------------------------------------------------------------------------
% LIMIETEN EN CONTINUÏTEIT. Alle opgaven zijn eigen werk.
%
% Opbouw: Start, leerstof (theorie, werkwijze of uitleg, voorbeeld, korte oefening), en op een
% nieuwe pagina de Opdrachten (door elkaar, stijgende moeilijkheid), Koppel en Test jezelf.
% Labels van dit hoofdstuk beginnen met h1-.
%----------------------------------------------------------------------------------------

\mychapter{1}{Limieten en continuïteit}

\lkr{Leerplan B+S'': LPD 26-30 (limieten en afgeleiden), het deel over het limietbegrip, rekenregels, onbepaalde vormen, asymptoten als limieten en continuïteit. De $\varepsilon$-$\delta$-definitie staat erin, met één uitgewerkt bewijs voor een eerstegraadsfunctie; de rest van het hoofdstuk rekent met de rekenregels. In Rationale functies werden asymptoten informeel ingevoerd: hier krijgen ze een definitie.}

Wat gebeurt er met $\frac{x^2 - 1}{x - 1}$ als $x$ naar $1$ gaat? Invullen kan niet: dat geeft $\frac00$. Toch komen de functiewaarden steeds dichter bij $2$. Een limiet beschrijft precies dat: waar de functiewaarden naartoe gaan, ook als de functie er zelf niet bestaat. Limieten zijn de basis van de afgeleide, die in hoofdstuk 2 volgt.

\begin{oefeningen}[Start]
\oef{h1-start-1}

\oef{h1-verken}
\end{oefeningen}

%========================================================================================
\section{Het begrip limiet}

\naast{%
Bekijk $f(x) = \frac{x^2 - 1}{x - 1}$ voor $x$ dicht bij $1$:

\vspace{1mm}
{\renewcommand{\arraystretch}{1.2}%
\begin{tabular}{c|cccc}
$x$ & $0{,}9$ & $0{,}99$ & $1{,}01$ & $1{,}1$ \\ \hline
$f(x)$ & $1{,}9$ & $1{,}99$ & $2{,}01$ & $2{,}1$
\end{tabular}}

\vspace{1mm}
Voor $x \neq 1$ is $f(x) = x + 1$: de grafiek is een rechte met een gat in $(1, 2)$. We schrijven $\lim_{x \to 1} f(x) = 2$.}{%
\begin{tikzpicture}[x=.6cm,y=.6cm,font=\scriptsize]
  \rooster{-1}{1}{3}{-1}{1}{4}
  \draw[accent,thick] (-1,0) -- (3,4);
  \open{(1,2)}
  \draw[grijs,dashed] (1,0) -- (1,2) (0,2) -- (1,2);
\end{tikzpicture}}

\begin{theorie}
$\formule{\lim_{x \to a} f(x) = L}$ betekent: $f(x)$ komt willekeurig dicht bij $L$ als $x$ voldoende dicht bij $a$ komt, met $x \neq a$. Nauwkeurig: voor elke $\varepsilon > 0$ bestaat er een $\delta > 0$ zodat
\[\formule{0 < |x - a| < \delta \Rightarrow |f(x) - L| < \varepsilon}.\]
De \textbf{\textcolor{accent}{linkerlimiet}} $\liml{x}{a} f(x)$ en de \textbf{\textcolor{accent}{rechterlimiet}} $\limr{x}{a} f(x)$ gebruiken alleen $x < a$, of alleen $x > a$. De limiet bestaat als beide bestaan en gelijk zijn.
\end{theorie}

De waarde $f(a)$ speelt geen rol: de functie hoeft in $a$ niet eens te bestaan. In de definitie is $\varepsilon$ de gevraagde nauwkeurigheid en $\delta$ hoe dicht $x$ bij $a$ moet komen om die te halen.

\voorbeeld{}{Bewijs met de definitie: $\lim_{x \to 2} (3x - 1) = 5$.\label{vb:h1-definitie}}{%
\uitleg{Neem een willekeurige $\varepsilon > 0$. Schrijf $|f(x) - L|$ uit en zoek een $\delta$ waarmee het kleiner wordt dan $\varepsilon$.}
\begin{stappenbreed}
|(3x - 1) - 5| = |3x - 6| = 3|x - 2| & Uitwerken. \\
3|x - 2| < \varepsilon \iff |x - 2| < \tfrac{\varepsilon}{3} & Wat er nodig is. \\
\delta = \tfrac{\varepsilon}{3}\colon\ 0 < |x - 2| < \delta \Rightarrow |f(x) - 5| < 3\delta = \varepsilon & Dat $\delta$ werkt.
\end{stappenbreed}}

\begin{oefeningen}
\oef{h1-begrip}
\end{oefeningen}

%========================================================================================
\section{Limieten berekenen}

\begin{theorie}
\begin{stelling}{Rekenregels}
Als $\lim_{x \to a} f(x) = L$ en $\lim_{x \to a} g(x) = M$, dan is
\[\formule{\lim_{x \to a} \bigl(f(x) \pm g(x)\bigr) = L \pm M} \qquad \formule{\lim_{x \to a} f(x)g(x) = LM} \qquad \formule{\lim_{x \to a} \frac{f(x)}{g(x)} = \frac{L}{M}} \ (M \neq 0).\]
\end{stelling}
Voor veeltermen, rationale functies (in hun domein), wortels, $\sin$, $\cos$, $\e^x$ en $\ln$ is de limiet in een punt van het domein de functiewaarde: je mag \textbf{\textcolor{accent}{invullen}}.
\end{theorie}

Geeft invullen $\frac00$, dan is dat een \textbf{onbepaalde vorm}: de limiet kan alles zijn. Vereenvoudig dan eerst.

\voorbeeld{}{Bereken $\displaystyle\lim_{x \to 3} \dfrac{x^2 - 9}{x - 3}$.\label{vb:h1-ontbinden}}{%
\uitleg{Invullen geeft $\frac00$. Ontbind en schrap de factor $x - 3$; dat mag, want $x \neq 3$.}
\begin{stappenbreed}
\lim_{x \to 3} \frac{(x - 3)(x + 3)}{x - 3} = \lim_{x \to 3} (x + 3) = 6 & Ontbinden, schrappen, invullen.
\end{stappenbreed}}

\voorbeeld{}{Bereken $\displaystyle\lim_{x \to 0} \dfrac{\sqrt{x + 4} - 2}{x}$.\label{vb:h1-wortel}}{%
\uitleg{Invullen geeft $\frac00$. Vermenigvuldig teller en noemer met de toegevoegde $\sqrt{x + 4} + 2$: dan verdwijnt de wortel uit de teller.}
\begin{stappenbreed}
\frac{(\sqrt{x + 4} - 2)(\sqrt{x + 4} + 2)}{x(\sqrt{x + 4} + 2)} = \frac{x + 4 - 4}{x(\sqrt{x + 4} + 2)} & $(a - b)(a + b) = a^2 - b^2$. \\
= \frac{1}{\sqrt{x + 4} + 2} \to \frac{1}{2 + 2} = \frac14 & Schrappen, invullen.
\end{stappenbreed}}

\begin{werkwijze}
Vul in. Krijg je een getal (met een noemer verschillend van nul), dan is dat de limiet. Krijg je $\frac00$: ontbind teller en noemer en schrap de gemeenschappelijke factor, of vermenigvuldig met de toegevoegde bij een wortel. Vul dan opnieuw in.
\end{werkwijze}

\begin{oefeningen}
\oef{h1-berekenen}
\end{oefeningen}

%========================================================================================
\section{Oneindige limieten en asymptoten}

Bij $\frac{1}{x - 1}$ dicht bij $1$ worden de functiewaarden onbeperkt groot: rechts van $1$ positief, links negatief. Voor grote $x$ naderen ze $0$. Dat schrijven we ook met limieten.

\begin{theorie}
$\formule{\lim_{x \to a} f(x) = +\infty}$ betekent: $f(x)$ wordt groter dan elk getal als $x$ voldoende dicht bij $a$ komt. Net zo voor $-\infty$ en voor eenzijdige limieten.

$\formule{\lim_{x \to +\infty} f(x) = L}$ betekent: $f(x)$ komt willekeurig dicht bij $L$ als $x$ voldoende groot wordt.

De rechte $x = a$ is een \textbf{\textcolor{accent}{verticale asymptoot}} als een eenzijdige limiet in $a$ gelijk is aan $\pm\infty$. De rechte $y = L$ is een \textbf{\textcolor{accent}{horizontale asymptoot}} als $\lim_{x \to +\infty} f(x) = L$ of $\lim_{x \to -\infty} f(x) = L$.
\end{theorie}

\voorbeeld{}{Bereken $\liml{x}{2} \dfrac{x + 1}{x - 2}$ en $\limr{x}{2} \dfrac{x + 1}{x - 2}$.\label{vb:h1-oneindig}}{%
\uitleg{De teller nadert $3$, de noemer $0$. Het teken van de noemer bepaalt het teken van het resultaat.}
\begin{stappenbreed}
\liml{x}{2} \frac{x + 1}{x - 2} = -\infty & Teller $\to 3$, noemer $\to 0$, negatief. \\
\limr{x}{2} \frac{x + 1}{x - 2} = +\infty & Noemer $\to 0$, positief.
\end{stappenbreed}}

\voorbeeld{}{Bereken $\displaystyle\lim_{x \to +\infty} \dfrac{3x^2 - x}{2x^2 + 5}$.\label{vb:h1-oneindig2}}{%
\uitleg{Invullen geeft $\frac{\infty}{\infty}$, ook een onbepaalde vorm. Deel teller en noemer door de hoogste macht van de noemer.}
\begin{stappenbreed}
\frac{3x^2 - x}{2x^2 + 5} = \frac{3 - \frac1x}{2 + \frac{5}{x^2}} & Delen door $x^2$. \\
\to \frac{3 - 0}{2 + 0} = \frac32 & $\frac1x \to 0$ en $\frac{5}{x^2} \to 0$.
\end{stappenbreed}
De grafiek heeft dus de horizontale asymptoot $y = \frac32$.}

\begin{oefeningen}
\oef{h1-oneindig}
\end{oefeningen}

%========================================================================================
\section{Continuïteit}

\begin{theorie}
Een functie $f$ is \textbf{\textcolor{accent}{continu}} in $a$ als $\formule{\lim_{x \to a} f(x) = f(a)}$: de limiet bestaat, $f(a)$ bestaat, en ze zijn gelijk. Ze is continu op een interval als ze continu is in elk punt ervan.

\begin{stelling}{Tussenwaardestelling}
Is $f$ continu op $[a, b]$ en hebben $f(a)$ en $f(b)$ een verschillend teken, dan heeft $f$ een nulwaarde in $\left]a, b\right[$.
\end{stelling}
\end{theorie}

Een continue functie op een interval teken je zonder je pen op te heffen. Veeltermen, $\sin$, $\cos$ en $\e^x$ zijn overal continu; rationale functies, wortels en $\ln$ op hun domein. Een functie die met verschillende voorschriften is gedefinieerd, kan op de overgang een sprong maken.

\voorbeeld{}{Bepaal $k$ zodat $f$ continu is: $f(x) = x^2$ voor $x \leqslant 1$ en $f(x) = 2x + k$ voor $x > 1$.\label{vb:h1-continu}}{%
\uitleg{Links en rechts van $1$ is $f$ een veelterm, dus continu. Alleen in $1$ moeten de linker- en rechterlimiet gelijk zijn aan $f(1)$.}
\begin{stappenbreed}
\liml{x}{1} f(x) = 1 = f(1),\quad \limr{x}{1} f(x) = 2 + k & Beide voorschriften invullen. \\
2 + k = 1 \iff k = -1 & Gelijkstellen.
\end{stappenbreed}}

\voorbeeld{}{Toon aan dat $x^3 + x - 1 = 0$ een oplossing heeft tussen $0$ en $1$. Zoek ze op één decimaal.\label{vb:h1-tussenwaarde}}{%
\uitleg{$f(x) = x^3 + x - 1$ is een veelterm, dus continu. Zoek een tekenwissel, en halveer het interval telkens.}
\begin{stappenbreed}
f(0) = -1 < 0,\ f(1) = 1 > 0 & Tussenwaardestelling: nulwaarde in $\left]0, 1\right[$. \\
f(0{,}5) = -0{,}375,\ f(0{,}75) \approx 0{,}17 & Tussen $0{,}5$ en $0{,}75$. \\
f(0{,}65) \approx -0{,}08,\ f(0{,}7) = 0{,}043 & Tussen $0{,}65$ en $0{,}7$: ongeveer $0{,}7$.
\end{stappenbreed}}

\begin{oefeningen}
\oef{h1-continu}
\end{oefeningen}

%========================================================================================
\opdrachtenkop

\begin{oefeningen}
\oef{h1-o1}

\oef{h1-o2}

\oef{h1-o3}

\oef{h1-o4}

\oef{h1-o5}

\oef{h1-o6}

\oef{h1-o7}

\oef{h1-o8}

\oef{h1-o9}
\end{oefeningen}

%----------------------------------------------------------------------------------------
% KOPPEL: verbanden met andere leerstof
%----------------------------------------------------------------------------------------
\needspace{20\baselineskip}
\opdrachtsectie{Koppel}

\lkr{Rationale functies, Rijen en Exponentiële en logaritmische functies: kijk na wat al gezien is. Het pakket legt de volgorde niet vast.}

\begin{oefeningen}
\oef{h1-koppel-1}

\oef{h1-koppel-2}

\oef{h1-koppel-3}
\end{oefeningen}

%----------------------------------------------------------------------------------------
\opdrachtsectie{Test jezelf}

\begin{oefeningen}
\oef{h1-t1}

\oef{h1-t2}

\oef{h1-t3}

\oef{h1-t4}

\oef{h1-t5}
\end{oefeningen}

\herhaalschema{%
  \fv{oef:h1-t1}{vb:h1-ontbinden}{oef:h1-berekenen}%
  \fv{oef:h1-t2}{vb:h1-wortel}{oef:h1-berekenen}%
  \fv{oef:h1-t3}{vb:h1-oneindig}{oef:h1-oneindig}%
  \fv{oef:h1-t4}{vb:h1-oneindig2}{oef:h1-oneindig}%
  \fv{oef:h1-t5}{vb:h1-continu}{oef:h1-continu}}

\ifzelfstudie
\section*{Maak je eigen samenvatting}
Sluit dit boek en schrijf in je schrift: wat $\lim_{x \to a} f(x) = L$ betekent (in woorden en met $\varepsilon$ en $\delta$), linker- en rechterlimiet, de rekenregels, wat je doet bij $\frac00$ en bij $\frac{\infty}{\infty}$, de definities van verticale en horizontale asymptoot, wat continu betekent, en de tussenwaardestelling. Vergelijk daarna met de kaders.
\fi
